\inicializarSeccion{Acomodo de cargas puntuales}
\tab Para el desarrollo de nuestra investigación de modelos de dipolos eléctricos, proponemos utilizar conjuntos de cargas eléctricas las cuales serán distribuidas en líneas rectas paralelas, por lo tanto se tomó la decisión de usar una distribución en líneas rectas horizontales, el motivo por el que decidimos utilizar este modelo fue porque conserva la esencia del modelo base donde una carga positiva se ubicaba a la derecha y una carga negativa se ubicaba del lado izquierdo.

\tab La carga positiva se encontrará ubicada en el punto $y = \frac{d}{2}$ y la carga negativa estará ubicada en el punto $y = -\frac{d}{2}$. Al estar acomodadas de esta manera, ambas líneas seran paralelas al eje de las “x” y estarán separadas por la distancia que dictamine a “d”, la carga positiva y la negativa estarán distribuidas por eje horizontal, así que siempre tendrán el mismo valor en $y$ pero el valor en $x$ estará variando dependiendo de la posición.